\section{A bijection with triangular-part compositions}\label{sec:coding}
Put $T_k=k(k+1)/2$ for $k\ge1$, and let $\Tri_N$ consist of the nonempty
compositions of $N$ whose parts belong to $\{T_1,T_2,\ldots\}$. This
class, its count and its reset/increment representations are known
\citep{wilson1998triangular,robbins2014polygonal}. We identify it with
the image of $F$; the test, identification and counting consequence form
one structural result.

\begin{theorem}[Explicit image coding]\label{thm:coding}
There is a weight-preserving bijection $\Phi:\Img_N\to\Tri_N$ given by
the following scan and inverse parsing.
\end{theorem}
\begin{proof}
Read a successful target from right to left. Its rightmost part $b\ge1$
initializes threshold $1$. From threshold $r$, a part $r+1$ increments the
threshold and a part at least $r+2$ resets it to $1$. Consequently, after
the rightmost part, there is a unique sequence of complete reset cycles
\begin{equation}\label{eq:cycle}
 2,3,\ldots,k,\;k+2+u\qquad(k\ge1,\ u\ge0)
\end{equation}
followed by a unique terminal increment string $2,3,\ldots,k$ with
$k\ge1$. A string is empty when $k=1$.

Start the encoded list with $b-1$ copies of $1$. For each complete cycle
\eqref{eq:cycle}, in \emph{scan order}, append $T_{k+1}$ followed by $u$
copies of $1$. Finally append $T_k$ for the terminal increment string.
The list is nonempty and has triangular parts. A complete cycle has mass
$(T_k-1)+(k+2+u)=T_{k+1}+u$, and the rightmost part together with the
terminal string has mass $b+T_k-1=(b-1)+T_k$. Thus total mass is preserved.

For the inverse, reserve the last triangular part, uniquely $T_k$, for
the terminal increment string. Parse the preceding list as an initial
string of ones and then parts at least $3$, each followed by its maximal
string of ones. This parse is unique since $1$ is the only positive
triangular number smaller than $3$. If there are $b-1$ initial ones,
begin a target read-list with $b$. Each subsequent triangular part is
uniquely $T_{j+1}$, $j\ge1$; when followed by $u$ ones, append
$2,3,\ldots,j,j+2+u$ to the read-list. Append $2,3,\ldots,k$ at the end,
then reverse the read-list to obtain the target.

Every reconstructed cycle and terminal string passes
\eqref{eq:scan}, so the inverse target belongs to $\Img_N$.
Its unique scan recovers the parsed pieces. Conversely, reserving the
last triangular part recovers the terminal string, even if that part is
$T_1=1$. These operations are mutually inverse, proving the bijection.
\end{proof}

For instance, $(2,3,2)$ encodes as $(1,3,3)$: the rightmost $2$ supplies
one initial $1$, the next $3$ is a reset, and the last $2$ is a terminal
increment. A one-part target $(N)$ encodes as $N$ ones.

\begin{corollary}[Known count, new image identification]\label{cor:series}
For $i_N=\lvert\Img_N\rvert$ and
$\Theta(z)=\sum_{k\ge1}z^{T_k}$, one has
\[
 \sum_{N\ge1}i_Nz^N=\frac{\Theta(z)}{1-\Theta(z)}.
\]
With the auxiliary convention $i_0=1$, this gives
$i_N=\sum_{T_k\le N}i_{N-T_k}$ for $N\ge1$.
\end{corollary}
\begin{proof}
Theorem~\ref{thm:coding} identifies each coefficient with the known
triangular-part count. A list of exactly $q\ge1$ triangular parts has
series $\Theta^q$, so summing gives the displayed identity. Taking off
the first part gives the recurrence. These are formal series: every
part has positive weight, hence every coefficient is finite and
$\Theta(0)=0$. The auxiliary empty list does not enlarge the carrier.
\end{proof}
